\documentclass[journal,10pt]{IEEEtran}

\usepackage{amsmath}
\usepackage{booktabs}
\usepackage{cite}
\usepackage{graphicx}
\usepackage[hidelinks]{hyperref}

\title{A Concise IEEEtran Article Template}

\author{First~A.~Author,~\IEEEmembership{Member,~IEEE,}
        Second~B.~Author,~\IEEEmembership{Senior~Member,~IEEE}
\thanks{Manuscript received Month Day, Year; revised Month Day, Year.}
\thanks{First A. Author and Second B. Author are with the Department of Example, Example University, City, Country (e-mail: author@example.edu).}}

\markboth{IEEE Transactions on Example,~Vol.~1, No.~1, Month~2026}
{Author \MakeLowercase{\textit{et al.}}: A Concise IEEEtran Article Template}

\begin{document}

\maketitle

\begin{abstract}
This template provides a compact IEEEtran starting point with the structural
elements most journal articles need: an abstract, index terms, sections,
equations, tables, figures, and numeric references.
\end{abstract}

\begin{IEEEkeywords}
IEEEtran, journal article, two-column format, LaTeX template.
\end{IEEEkeywords}

\section{Introduction}
\IEEEPARstart{T}{his} opening paragraph uses the IEEE drop-cap convention.
Summarize the problem, the gap in prior work, and the main contribution.

\section{Method}
Describe the proposed method in reproducible terms. Use equations where they
clarify the core computation:
\begin{equation}
  y = f_{\theta}(x) + \epsilon.
\end{equation}

\section{Results}
Table~\ref{tab:results} shows a compact comparison using booktabs.

\begin{table}[!t]
\caption{Example Quantitative Results}
\label{tab:results}
\centering
\begin{tabular}{lcc}
\toprule
Method & Accuracy & F1 \\
\midrule
Baseline & 0.812 & 0.798 \\
Proposed & 0.874 & 0.861 \\
\bottomrule
\end{tabular}
\end{table}

\begin{figure}[!t]
\centering
\fbox{\rule{0pt}{1.1in}\rule{0.9\columnwidth}{0pt}}
\caption{Replace this placeholder with a publication-ready figure.}
\label{fig:placeholder}
\end{figure}

\section{Conclusion}
Restate the main result, note limitations, and identify the next concrete
research step.

\begin{thebibliography}{1}

\bibitem{ieeetran}
M.~Shell, ``How to use the IEEEtran LaTeX class,'' \emph{Journal of LaTeX Class Files}, vol.~14, no.~8, pp.~1--8, 2015.

\end{thebibliography}

\end{document}
